\documentclass[12pt,onecolumn,letterpaper]{article}
\usepackage[spanish,es-nodecimaldot]{babel}
\usepackage[utf8]{inputenc}
\usepackage[T1]{fontenc}
\usepackage{float}
\usepackage{graphicx}
\usepackage{microtype}

%% Sets page size and margins
\usepackage[a4paper,top=2.5cm,bottom=2.5cm,left=2cm,right=2cm]{geometry}

\usepackage{amsmath}
\usepackage{amssymb}
\usepackage{amsthm}
\usepackage{mathtools}
\usepackage{listings}
\usepackage{enumitem}[shortlabels]
\usepackage[colorinlistoftodos]{todonotes}
\usepackage[colorlinks=true, allcolors=blue]{hyperref}
\usepackage{bookmark}
\usepackage{url}
\usepackage{lipsum}
\usepackage{tcolorbox}
\usepackage{tikz}

\usepackage{listings}
\lstset{
    language=bash,
    basicstyle=\ttfamily
}

%Author affil
\usepackage{authblk}

%Para la parte de las imagenes que estarán en assets 
\graphicspath{{assets/}}

%% Title
\title{
		\vspace{-0.7in} 	
		\usefont{OT1}{bch}{b}{n}
		
		\begin{minipage}{3cm}
            \vspace{-0.5in} 	
    	    \begin{center}
    		\includegraphics[height=3.2cm]{unam.png} %Parte donde va el logo de la unam 
    	\end{center}
    \end{minipage}\hfill
    \begin{minipage}{10.7cm}
    	\begin{center}
            \normalfont \normalsize \textsc{UNIVERSIDAD NACIONAL AUTÓNOMA DE MÉXICO \\ FACULTAD DE CIENCIAS \\ Modelado y Programación} \\
		    \huge Proyecto 2 \\ Creación de un Reproductor de Música con una Base da Datos SQL\\
        
            \normalsize \rmfamily \mdseries
            \vspace{0.1in}
            Sánchez Girón Yael Adrian$^1$\\

        \vspace{0.1in}
        $^1$Facultad de Ciencias, UNAM\\
        \date{\date{\vspace{0.1in}\today}}
    \end{center}
    
\end{minipage}\hfill
\begin{minipage}{3.2cm}
    \vspace{-0.5in} 
    \begin{center}
    	\includegraphics[height=3.2cm]{fcunam.png}
    \end{center}
\end{minipage}

\author{}

\vspace{0.1in}
\begin{tcolorbox}[colframe=white, arc=8pt]
    \normalsize \rmfamily \mdseries
    \centering
    \textbf{Resumen}\\

% Puedes colocar el texto del resumen en esta parte
Este reporte presenta el desarrollo de un sistema de reproducción de música con una base relacional de datos SQLite, utilizando los lenguajes de Rust y Flutter, que para operar en conjunto se utilizó la herramienta de Flutter\_Rust\_Bridge, para poder integrar ambos lenguajes, se implementaron los patrones de DAO y MVC para el correcto desglose de responsabilidades del programa, además de poder desarrollar la aplicación para los sistemas operativos de macOS y Linux.
Se llegó a un sistema que cumple con los requerimientos del proyecto, es decir, procesar archivos ID3V2.4, tener una base de datos la cual almacene los datos de las canciones (Performer, Nombre, Grupo, etc), una interfaz gráfica la cual es capaz de poder reproducir los audios, y una barra de búsqueda. 

    \medskip
\end{tcolorbox}
\vspace{-0.7in} 
}




\hypersetup{
    colorlinks=true,
    linkcolor=blue,
    filecolor=magenta,      
    urlcolor=cyan
    }


%------Vectores Unitarios-------%
\newcommand{\uveci}{\hat{\textbf{i}}} 
\newcommand{\uvecj}{\hat{\textbf{j}}}
\newcommand{\uveck}{\hat{\textbf{k}}}

\newcommand{\grade}{^{\circ}}
    \setlength {\marginparwidth }{2cm}

\usepackage{listings}
\lstset{
    basicstyle=\ttfamily\small, % Fuente de código 
    breaklines=true,            % Activa el salto de línea automático
    columns=fullflexible
}

\begin{document}

\maketitle

\section{Introducción}
Para la realización de este segundo proyecto de la materia de Modelado y Programación se nos encomendó la realización de un Reproductor de Música, el cual implementa una base de datos SQLite, misma que se nos fue proporcionada al inicio del mismo proyecto.
Las especificaciones de este programa, fueron varias, entre los que más destacan son:
\begin{enumerate}
    \item Que el reproductor sea capaz de leer las entradas de la base de datos que se nos fue proporcionada
    \item Los archivos de audio deben de tener la etiqueta de ID3V2.4, y con los metadatos que contenga el archivo se llenen los parámetros establecidos en la base de datos
    \item El archivo .db que se genere para la base de datos se almacene dentro del directorio estándar de datos de cada sistema operativo.
    \item A diferencia del primer proyecto, en donde era opcional la implementación de una interfaz gráfica, en este proyecto se exige la implementación de una.
    \item Dentro de la interfaz gráfica debe de ser capaz de poder editar las entradas de la base de datos, cabe aclarar que se especifica que los archivos de audio NO deben de ser modificados.
    \item La implementación de un ``minero'' el cual se encargará de escanear los archivos mp3, que sean ID3V2, y que además los inserte dentro de la base de datos, para que así, los pueda reproducir la interfaz gráfica
    \item Que se puedan reproducir los archivos dentro de la interfaz gráfica
    \item Cada uno de los componentes (Minero, Reproductor) deben de tener su propio ejecutable
\end{enumerate}
Con estos parámetros fue como se inicio con la parte del análisis y modelado del programa.
\section{Modelado y Programación}
% Recordar que esta es la parte más larga de todo el documento

La primera cosa se tuvo a consideración para la realización de este proyecto fue, con qué lenguajes de programación se tendría que trabajar, se optó con la implementación de los mismos dos lenguajes de programación utilizados en el proyecto pasado, es decir, se utilizaron los lenguajes de \textbf{Rust} y el toolkit de \textbf{Flutter} (que utiliza el lenguaje de programación \textbf{Dart}), esto ya que existe cierta familiaridad, tanto con la sintaxis, como los paradigmas que implementan ambos lenguajes de programación, esto porque ya fueron implementados anteriormente.

Antes de saber que bibliotecas utilizar en cada lenguaje de programación se realizó el desglose de responsabilidades, de tal forma que se pudo dividir en 3 sub-categorías, siendo estas:
\begin{enumerate}
    \item Minero
    \item DAO (\textit{Data Access Object})
    \item Reproductor (\textit{Interfaz Gráfica})
\end{enumerate}
%Colocar la cita de Flutter Rust Bridge
Con base en esta división, se estructuró un espacio de trabajo (\textit{workspace}) compuesto por tres sub-proyectos dentro del repositorio principal: el DAO y el minero desarrollados en Rust, y la interfaz gráfica desarrollada en Flutter/Dart. Para comunicar ambas capas se integró \texttt{flutter\_rust\_bridge}, herramienta que genera los enlaces necesarios para invocar funciones de Rust de manera nativa y segura desde Dart. A continuación se muestra la estructura de directorios resultante:
%Poner el diagrama de el directorio principal con los tres subdirectorios que son los sub-proyectos descritos

A partir de aquí procedemos con el análisis de cada uno de los ``sub-proyectos'' que constituyen en su totalidad el reproductor.

\subsection{DAO}
Para la parte del DAO(Data Access Object), en inicio este término hace referencia al patrón de diseño, pero en términos del proyecto se opto por nombrar a este ``Sub-proyecto'' ya que la tarea que se debe de realizar es la de, tener las funciones necesarias para poder: Insertar, Modificar, Eliminar cada una de las entradas de la base de datos, esto para que no existan dichas funciones dentro del minero y de la interfaz gráfica y que solamente se manden a llamar en caso de ser requeridas, en esta parte se optó por que cada función de agregar uno de los parámetros de cada tabla, se devolviera el número generado de cada nuevo ID, y en las funciones que fueran de editar o de eliminar, cada una de las funciones debe de devolver un booleano, para confirmar que la acción realizada fue exitosa o hubo algún fallo. Esta parte del DAO, fue la más simple de las 3 ya que solo fue la implementación de las funciones descritas anteriormente, aunque en etapas más avanzadas del proyecto se tuvo que editar y agregar nuevas funciones que cumplieran con las necesidades que requería el proyecto en ese momento, cabe mencionar que en este archivo es donde se implementó la función para que el archivo .db se generara en la ruta de acuerdo con las especificaciones de cada sistema operativo, por ejemplo, en Linux se genera en la ruta \textit{.local/share} o en macOS se genera en la ruta \textit{Library/Application Support}.

Además, que para los comandos que se requirieron dentro del SQL se tuvo que consultar que hacían los comandos SQL más utilizados, de los cuales se encuentran comandos como \textbf{SELECT}, \textbf{INSERT INTO}, \textbf{DELETE}, entre otros muchos, los cuales fueron investigados de una página de internet llamada W3Schools %insertar la fuente 

\subsection{Minero}
Para el Minero (en Inglés conocido como un Metadata Miner), se tuvieron los siguientes objetivos:
\begin{enumerate}
    \item Poder recorrer de forma recursiva el directorio y filtrar los archivos que tengan la etiqueta ID3V2.4 (File Crawler)
    \item Poder extraer los metadatos de cada archivo que contenga las etiquetas antes mencionadas (ID3 Parser)
    \item El Minero el cual es el que utiliza el File Crawler y el Parser para poder realizar sus funciones además de implementar las funciones del DAO para poder insertar las entradas en la base de datos
    \item Y un visualizer para poder ejecutar el minero de forma independiente
\end{enumerate}

Una vez teniendo los objetivos esclarecidos, se inicio con la búsqueda de las bibliotecas que nos ayudaron con la creación del minero, de las cuales destacan varias, entre ellas se encuentran las bibliotecas de \textbf{dirs}, \textbf{walkdir}, \textbf{id3}, \textbf{rusqlite}.
%Poner las referencias de los crates 

Gracias a estas bibliotecas la implementación del minero resultó ser de las más sencillas para la parte de la implementación del minero, ya que pues la mayoría si no es que todas las funciones necesarias para poder armar el minero. Para poder hacer pruebas reales del minero se optó por descargar música .mp3 desde YouTube, para esto se hizo uso de la herramienta \textbf{yt-dlp} la cual permitió descargar los archivos con las etiquetas y datos necesarios para comprobar que el minero funcionara como debería. %Poner referencia y una imagen del minero funcionando

\subsection{Interfaz Gráfica}
Esta fue la parte más extensa de todo el proyecto, esto debido a que hubo varios errores en la toma de excepciones, y también se consumió mucho tiempo en el diseño de cómo debería quedar la interfaz gráfica, al final se opto por las siguientes características: 
\begin{enumerate}
    \item Que el centro de la aplicación tenga dos vistas principales, una de Grid, donde se pueda ver de forma general las portadas de los tracks, y otra vista en forma de lista donde además se tendrá un botón para poder editar las canciones
    \item De ahí en esta última vista, el botón de editar debe de mostrar un pop-up en la parte inferior que se debe de indicar que se va a editar, el track o el performer, de aquí se tiene que especificar si el performer es una persona o un grupo, y dependiendo de la opción que se escoja, se edita con los parámetros que se requieran para dicho campo, dentro de la pestaña de editar el grupo, se tiene otro botón el cual permite el poder relacionar a las personas con los grupos, de esta forma, podemos buscar a una persona que está en cierto grupo y viceversa.
    \item En la parte inferior se encuentra el reproductor, donde, en la parte izquierda aparece la portada del álbum de forma ``miniatura'', seguido de los botones, de play/pausa, avanzar, retroceder y una barra de reproducción para poder manejar el tiempo a disposición del usuario, y finalmente en la parte derecha un control deslizante para el manejo del volumen
    \item En la parte superior se encontrará la barra de búsqueda la cual podrá filtrar a través de los parámetros de: 
    \begin{enumerate}
        \item Título
        \item Artista
        \item Álbum
        \item Género
        \item Año (year)
    \end{enumerate}
    los cuáles se encontraran en un botón en la derecha de la barra de búsqueda, además de un botón de poder borrar toda  la búsqueda realizada
\end{enumerate}

Para poder entender de que forma esta parte de la interfaz gráfica se integrará a los demás componentes, se tuvo que realizar un diagrama de modelo vista controlador, donde los ``participantes'' se pueden dividir como: 
\begin{enumerate}
    \item \textbf{Modelo:} Comprende las estructuras de (\texttt{Rola, RolaBD, MetadatosCancion}), la capa de la base de datos en SQLite, la cual es manejada por el DAO, el minero que extrae la metadata de los archivos con etiqueta ID3V2.4 y el motor de reproducción de Audio
    \item \textbf{Controlador:} Agrupa las clases ``reactivas'' basadas en \texttt{ChangeNotifier}, aquí se encuentran los controladores de búsqueda, reproductor, biblioteca, encargadas de coordinar las llamadas hacia Rust, mantener el estado en memoria y notificar los cambios a la interfaz
    \item \textbf{Vista:} Conformada por todos los \textit{widgets} y pantallas \textit{screens} construidos en Flutter, los cuales tienen el objetivo de renderizar la interfaz gráfica, y capturar las interacciones del usuario
\end{enumerate}
%Poner en bonito el diagramota donde aparece Paco :D 

Cabe mencionar que, para la parte de la vista se optó por separar a cada una de las partes de la interfaz gráfica por cada una de sus funcionalidades (features), esto para tener dentro del proyecto los directorios de: 
\begin{enumerate}
    \item Biblioteca (Donde se encuentran las vistas de Grid y lista)
    \item Buscador (Con sus respectivo botón con la opción de filtro)
    \item Edición (Para poder seleccionar si editar el track o el artista/grupo y depende de cada opción tendrá sus propias características)
    \item Reproductor (La barra de reproducción)
    \item Layout Principal (Donde se juntaran los elementos de biblioteca, buscador, edición y reproductor)
\end{enumerate} 
Esto permitió una mejor separación de cada una de las responsabilidades de cada parte de la interfaz gráfica, aunque aún hay cosas que deben de tener una refactorización, por ejemplo, dentro del archivo de \textit{controlador\_biblioteca} el cual contiene la clase de Rola, el cual contiene de atributos todas las características que se pueden obtener de un archivo \textbf{ID3V2.4}, sin embargo esta clase se utiliza a lo largo de otros archivos dentro del reproductor, por lo que lo más sensato es poner dicha clase en un archivo a un nivel ``superior'' dentro de la jerarquía de los directorios, de esta forma, el resto de archivos puede tener acceso a dicha clase. 

\section{Implementación}

La materialización del diseño arquitectónico mencionado anteriormente se llevó a cabo dentro de un repositorio unificado bajo la estructura de un \textit{workspace} de Cargo, el cual integra los sub-proyectos de Rust (\texttt{minero} y \texttt{nucleo\_dao}), el esquema relacional (\texttt{base\_de\_datos}) y la aplicación de escritorio en Flutter (\texttt{reproductor}).

\subsection{Tecnologías utilizadas}
Para la construcción de cada uno de los sub-proyectos se seleccionaron las siguientes herramientas y bibliotecas: 

\begin{itemize}
    \item \textbf{Backend Rust}
    \begin{itemize}
        \item \texttt{rusqlite}: Controlador nativo para SQLite utilizado en \texttt{nucleo\_dao} para gestionar conexiones, transacciones ACID y restricciones de integridad referencial mediante \texttt{PRAGMA foreign\_keys = ON}.
        \item \texttt{id3}: Biblioteca encargada de analizar las etiquetas \texttt{ID3v2.4} de los archivos de audio, permitiendo extraer metadatos textuales y portadas incrustadas (\textit{CoverFront} y \textit{Artist}).
        \item \texttt{walkdir}: Utilizada en el módulo de exploración para recorrer directorios de forma recursiva y localizar archivos con extensión \texttt{.mp3} y cabecera válida \texttt{ID3}.
        \item \texttt{dirs}: Permite resolver de forma multiplataforma las rutas estándar de datos de usuario para alojar el archivo \texttt{rolas.db} sin depender de rutas relativas del proyecto.\end{itemize}

\item \textbf{Interoperabilidad (FFI):}
    \begin{itemize}
        \item \texttt{flutter\_rust\_bridge}: Herramienta de generación de enlaces (\textit{bindings}) que compila los \textit{crates} de Rust en una librería dinámica nativa (\texttt{rust\_lib\_reproductor}) y expone sus funciones de manera asíncrona y segura hacia Dart.\par
    \end{itemize}

\item \textbf{Interfaz gráfica, estado y reproducción (Flutter / Dart):}
    \begin{itemize}
        \item \texttt{macos\_ui}: Biblioteca principal de \textit{widgets} de escritorio que define la estética de ventanas, barras de herramientas y controles, complementada con componentes base de Flutter (\texttt{material} y \texttt{cupertino}).\cite{pubdev_cupertinotabview}
        \item \texttt{provider}: Gestor de estado reactivo basado en \texttt{ChangeNotifier}, empleado por los controladores de biblioteca, búsqueda, para poder actualizar la interfaz en tiempo real al filtrar los tracks o editar los datos de la base de datos SQLite \cite{flutter_widget_previewer}
        \item \texttt{just\_audio} (junto con \texttt{just\_audio\_media\_kit} y \texttt{media\_kit\_libs\_linux}): Motor encargado de la carga, decodificación y reproducción de los archivos \texttt{.mp3}, así como del control de progreso y volumen. En Linux se inicializa condicionalmente sobre \texttt{libmpv} para dotar de soporte nativo de audio al entorno de escritorio.\cite{pubdev_file_picker,pubdev_macos_ui_package,pubdev_just_audio,pubdev_media_kit_libs_linux}
        \item \texttt{file\_picker}: Componente multiplataforma utilizado para desplegar el explorador de archivos del sistema operativo y seleccionar los directorios de música a escanear.\cite{pubdev_file_picker}
        \item \texttt{shared\_preferences}: Almacenamiento ligero clave-valor para conservar preferencias de la aplicación entre ejecuciones.\cite{share}
    \end{itemize}
    %Utilizar las referencias de su documentación oficial en cada uno de los puntos establecidos

\subsubsection{Requerimientos del Sistema}

El proyecto contempla dos modalidades de uso: la ejecución directa mediante los binarios publicados en la sección \textit{Releases} del repositorio, o la compilación manual a partir del código fuente.

\begin{enumerate}
    \item \textbf{Para ejecutar los binarios precompilados:}
    \begin{itemize}
        \item Este proyecto está dirigido a sistemas operativos \textbf{macOS} (binario universal compatible con arquitecturas Apple Silicon \texttt{arm64}  y \textbf{Linux} de 64 bits (\texttt{x86\_64}).
        \item El esquema relacional (\texttt{rolas.sql}) y las librerías nativas de reproducción (\texttt{libmpv} mediante \texttt{media\_kit\_libs\_linux}) se empaquetan dentro de la distribución, por lo que no es necesario instalar Flutter, Rust ni gestores externos de bases de datos.
        \item El archivo \texttt{rolas.db} se crea e inicializa automáticamente en la primera ejecución siguiendo las convenciones de cada sistema operativo: el estándar de Apple en macOS (\texttt{\textasciitilde/Library/Application Support/rolas/}) y la especificación XDG en Linux (\texttt{\textasciitilde/.local/share/rolas/}).
    \end{itemize}

    \item \textbf{Para compilar desde el código fuente:}
    \begin{itemize}
        \item \textit{Herramientas base:} Git, SDK de Flutter (con soporte de escritorio habilitado) y el \textit{toolchain} de Rust (\texttt{rustc}, \texttt{cargo}).
        \item \textit{Dependencias en macOS:} Xcode Command Line Tools y CocoaPods para enlazar las bibliotecas nativas de Rust.
        \item \textit{Dependencias de Linux:}
        \begin{enumerate}
            \item[a)] Compiladores y herramientas de construcción: \texttt{clang}, \texttt{cmake}, \texttt{ninja-build}, \texttt{pkg-config}, así como \texttt{g++} (Debian/Ubuntu) o \texttt{gcc-c++} (Fedora).
            \item[b)] Cabeceras de desarrollo de GTK: \texttt{libgtk-3-dev} (Debian/Ubuntu) o \\ \texttt{gtk3-devel} (Fedora).
            \item[c)] SQLite: \texttt{libsqlite3-dev} (Debian/Ubuntu) o \texttt{sqlite-devel} (Fedora).
            \item[d)] Salida y decodificación de audio: \texttt{libasound2-dev} y \texttt{libmpv-dev} \\ (Debian/Ubuntu), o \texttt{alsa-lib-devel} y \texttt{mpv-devel} (Fedora, habilitando el repositorio RPM Fusion Free).
            \item[e)] Diálogos nativos de selección de carpetas: \texttt{zenity}.
            \item[f)] Soporte de verificación para \texttt{cargokit}: \texttt{perl-Digest-SHA} en Fedora (requerido por \texttt{shasum} al compilar el puente de Rust cuando Flutter se instala mediante Snap).
        \end{enumerate}
    \end{itemize}
\end{enumerate}

\subsection{Ejecución del programa}
Ahora, se detallarán los comandos necesarios para poder clonar, probar, ejecutar y empaquetar el sistema desde la terminal: 

\begin{enumerate}
    \item \textbf{Clonar el repositorio:}
\begin{lstlisting}
git clone https://github.com/A4DREK/Segundo-Proyecto-Modelado-y-Programaci-n.git
cd Segundo-Proyecto-Modelado-y-Programaci-n
\end{lstlisting}
    \item \textbf{Ejecutar las pruebas unitarias del backend (DAO y Minero):}
\begin{lstlisting}
cargo test --workspace
\end{lstlisting}
    \item \textbf{Ejecutar el minero de forma independiente (CLI):}
    Permite escanear un directorio específico y mostrar en la CLI un log de que canciones aparecen.
\begin{lstlisting}
cargo run -p minero -- /ruta/a/carpeta/musica\end{lstlisting}
    \item \textbf{Ejecutar la interfaz gráfica en modo desarrollo:}
\begin{lstlisting}
cd reproductor
flutter pub get
flutter run
\end{lstlisting}
    \item \textbf{Generar los ejecutables independientes (\textit{Release}):}
\begin{lstlisting}
# En macOS (genera build/macos/Build/Products/Release/reproductor.app)
flutter build macos --release

# En Linux (genera el paquete autocontenido en build/linux/x64/release/bundle/)
flutter build linux --release
\end{lstlisting}

\end{enumerate}


\section{Conclusiones}

El desarrollo de este proyecto permitió el desarrollo de habilidades con relación a una introducción a los comandos de Bases de Datos SQL y una profundización del desarrollo del patrón Modelo-Vista-Controlador esto a través de las especificaciones dadas por el propio proyecto del reproductor de música.

La elección de tener una arquitectura ``híbrida'' con dos lenguajes de programación (Rust y Flutter) permitió el desglose de responsabilidades y además, fue una decisión importante a la hora de equilibrar el rendimiento y el escribir la interfaz de usuario, esto ya que Rust, se caracteriza por la seguridad en memoria, y rapidez, esto por su sistema de ``ownership'', el cual se explicó en el proyecto pasado, y por su parte Flutter, al ser un toolkit destinado a la creación de interfaces gráficas multiplataforma, resultó poder crear la interfaz de una forma un poco más clara, además de poder ser reactiva y modular, esto para poder llevar acabo de buena forma el patron Modelo-Vista-Controlador, y tener un modelo basado en funcionalidades como se describió anteriormente, mientras que utilizar la herramienta de Flutter\_Rust\_Bridge, mostró ser una de las opciones más óptimas para poder integrar estos dos lenguajes de programación, esto gracias a que es posible invocar funciones de Rust de forma asíncrona sin comprometer la fluidez de la interfaz gráfica. 

Asimismo, la separación del sistema en tres módulos independientes (\texttt{minero}, \texttt{nucleo\_dao} y \texttt{reproductor}) aportó ventajas clave en mantenimiento y pruebas:
\begin{itemize}
    \item El encapsulamiento de las consultas en el patrón \textbf{DAO} evitó la duplicación de código SQL y permitió validar cada operación de persistencia mediante pruebas unitarias automatizadas (\texttt{cargo test}) antes de integrarla con la interfaz.
    \item El diseño del \textbf{minero} como componente autónomo con salida tabular en consola (\textit{visualizer}) posibilitó depurar la extracción de metadatos de manera aislada y cumplir con el requisito de contar con ejecutables independientes.
    \item La edición de registros directamente sobre la base de datos SQLite aseguró el cumplimiento de la restricción no destructiva del sistema, permitiendo corregir información de canciones, artistas, grupos y relaciones sin alterar los archivos de audio originales.
    \item El uso de rutas estándar dependientes del sistema operativo y la incrustación del esquema SQL dentro del binario garantizaron la portabilidad de los ejecutables tanto en \textbf{macOS} como en \textbf{Linux}.
\end{itemize}

Finalmente, como trabajo futuro y oportunidades de mejora se identifica la refactorización de las clases de dominio compartidas (como la clase \texttt{Rola}) hacia una capa superior en Flutter para reducir el acoplamiento entre controladores, la ampliación del minero para soportar formatos adicionales sin pérdida (como \texttt{FLAC} o \texttt{WAV}) y la incorporación de listas de reproducción personalizadas persistentes en la base de datos.

\nocite{*}
\bibliographystyle{plain}
\bibliography{references}
\end{itemize}
\end{document}